% ============================================================
%  第二章：LaTeX 常用元素示例
%  本章展示公式、定理、表格、图片、交叉引用等常见用法。
%  实际写作时请将本章内容替换为自己的章节。
% ============================================================
\chapter{第二章标题}

% ────────────────────────────────────────────────────────────
\section{数学公式}
% ────────────────────────────────────────────────────────────

\subsection{行内公式与单行公式}

行内公式直接用美元符号包裹，例如 $f(x) = x^2 + 1$ 或
薛定谔方程 $\hat{H}\psi = E\psi$。

带编号的单行公式使用 \texttt{equation} 环境：
\begin{equation}
  \int_{-\infty}^{+\infty} e^{-x^2}\,\dd x = \sqrt{\pi}.
  \label{eq:gaussian}
\end{equation}

引用时用 \cref{eq:gaussian}（或 \eqref{eq:gaussian}）。

\subsection{多行对齐公式}

使用 \texttt{align} 环境在等号处对齐：
\begin{align}
  (a+b)^2 &= a^2 + 2ab + b^2, \label{eq:binom2} \\
  (a-b)^2 &= a^2 - 2ab + b^2. \label{eq:binom2m}
\end{align}

不需要编号的行在末尾加 \verb|\notag| 或改用 \texttt{align*}。

\subsection{分段函数}

\begin{equation}
  |x| =
  \begin{cases}
    \phantom{-}x, & x \geq 0, \\
    -x,           & x < 0.
  \end{cases}
\end{equation}

\subsection{矩阵}

\begin{equation}
  A = \begin{pmatrix} a_{11} & a_{12} \\ a_{21} & a_{22} \end{pmatrix},
  \qquad
  \det(A) = a_{11}a_{22} - a_{12}a_{21}.
\end{equation}

% ────────────────────────────────────────────────────────────
\section{定理、引理与证明}
% ────────────────────────────────────────────────────────────

% 定理环境在 main.tex 中已定义（theorem / lemma / proposition 等）

\begin{theorem}[勾股定理]\label{thm:pythagoras}
  设 $a$、$b$ 为直角三角形的两直角边，$c$ 为斜边，则
  \[
    a^2 + b^2 = c^2.
  \]
\end{theorem}

\begin{proof}
  （在此给出证明过程。）
\end{proof}

\begin{lemma}\label{lem:example}
  （引理的内容。）
\end{lemma}

\begin{corollary}
  由\cref{thm:pythagoras} 可知，……
\end{corollary}

% ────────────────────────────────────────────────────────────
\section{表格}
% ────────────────────────────────────────────────────────────

使用 \texttt{booktabs} 宏包（已在 \texttt{main.tex} 中加载）绘制三线表，
如\cref{tab:example}所示。

\begin{table}[htbp]
  \centering
  \caption{示例表格：不同方法在基准上的对比}
  \label{tab:example}
  \begin{tabular}{l c c c}
    \toprule
    方法           & 指标 1 & 指标 2 & 指标 3 \\
    \midrule
    方法 A         & 0.832  & 1.24   & 97.3\% \\
    方法 B         & 0.751  & 1.56   & 95.1\% \\
    \textbf{本文方法} & \textbf{0.891} & \textbf{1.10} & \textbf{98.6\%} \\
    \bottomrule
  \end{tabular}
\end{table}

% 若表格较宽，可改用 longtable 或旋转（rotating 宏包的 sidewaystable）。

% ────────────────────────────────────────────────────────────
\section{插图}
% ────────────────────────────────────────────────────────────

图片建议以 PDF 或 PNG 格式存放在 \texttt{figs/} 子目录中。
\Cref{fig:example} 展示了插图的基本用法。

\begin{figure}[htbp]
  \centering
  % 取消注释并替换文件名：
  % \includegraphics[width=0.6\textwidth]{figs/your_figure.pdf}
  \framebox[0.6\textwidth]{\rule{0pt}{4cm}\centering 图片占位}
  \caption{图片标题：说明图片内容。}
  \label{fig:example}
\end{figure}

\subsection{并排两图}

\begin{figure}[htbp]
  \centering
  % 左图
  \begin{minipage}[b]{0.45\textwidth}
    \centering
    % \includegraphics[width=\textwidth]{figs/fig_a.pdf}
    \framebox[\textwidth]{\rule{0pt}{3cm}\centering 左图}
    \caption{左图标题。}
    \label{fig:left}
  \end{minipage}
  \hfill
  % 右图
  \begin{minipage}[b]{0.45\textwidth}
    \centering
    % \includegraphics[width=\textwidth]{figs/fig_b.pdf}
    \framebox[\textwidth]{\rule{0pt}{3cm}\centering 右图}
    \caption{右图标题。}
    \label{fig:right}
  \end{minipage}
\end{figure}

% ────────────────────────────────────────────────────────────
\section{列表}
% ────────────────────────────────────────────────────────────

无序列表：
\begin{itemize}
  \item 第一条。
  \item 第二条。
  \item 第三条。
\end{itemize}

有序列表：
\begin{enumerate}
  \item 步骤一：……
  \item 步骤二：……
  \item 步骤三：……
\end{enumerate}

% ────────────────────────────────────────────────────────────
\section{文献引用与交叉引用}
% ────────────────────────────────────────────────────────────

引用文献用 \verb|\cite{key}|，例如见文献~\cite{doe2020}。
多篇同时引用：\cite{doe2020,smith2021}。

引用本文中的公式、定理、图、表分别用 \verb|\cref{}|：
\cref{eq:gaussian}、\cref{thm:pythagoras}、\cref{fig:example}、\cref{tab:example}。

% ────────────────────────────────────────────────────────────
\section{本章小结}
% ────────────────────────────────────────────────────────────

本章介绍了……（总结本章主要内容）。
